% Part: first-order-logic
% Chapter: completeness
% Section: completeness-theorem

\documentclass[../../../include/open-logic-section]{subfiles}

\begin{document}

\iftag{FOL}
      {\olfileid[pt]{fol}{com}{cth}}
      {\olfileid[pt]{pl}{com}{cth}}

\olsection{O Teorema da Completude}

\begin{explain}
Vamos combinar nossos resultados: chegamos ao teorema da completude.
\end{explain}

\begin{thm}[Teorema da Completude]
\ollabel{thm:completeness}
Seja $\Gamma$ um conjunto de !!{sentence}s.  Se $\Gamma$ é consistente, então
é satisfatível.
\end{thm}

\begin{proof}
Suponha que $\Gamma$ seja consistente.\iftag{FOL}{ Por
  \olref[hen]{lem:henkin}, existe um conjunto saturado consistente
  $\Gamma' \supseteq \Gamma$.}{} Por \olref[lin]{lem:lindenbaum}, existe
um $\Gamma^* \supseteq \iftag{FOL}{\Gamma'}{\Gamma}$ que é
consistente e !!{complete}.
\iftag{FOL}{
  Como $\Gamma' \subseteq \Gamma^*$, para cada
  !!{formula}~$!A(x)$, $\Gamma^*$ contém !!a{sentence} da forma
  \iftag{prvEx}
        {$\lexists[x][!A(x)] \lif !A(c)$}
        {$\lnot\lforall[x][!A(x)] \lif \lnot !A(c)$}
  e, assim, $\Gamma^*$ é saturado.  Se $\Gamma$
  não contém~$\eq$, então, por}{Por}
\olref[mod]{lem:truth},
$\iftag{FOL}{\Sat{M(\Gamma^*)}}{\pSat{v(\Gamma^*)}}{!A}$ se e somente se $!A \in
\Gamma^*$.  Disso segue, em particular, que, para todo $!A \in
\Gamma$, $\iftag{FOL}{\Sat{M(\Gamma^*)}}{\pSat{v(\Gamma^*)}}{!A}$, de modo que
$\Gamma$ é satisfatível.\iftag{FOL}{
  Se $\Gamma$ de fato contém~$\eq$,
  então, por \olref[ide]{lem:truth}, para todas as !!{sentence}s~$!A$,
  $\Sat{\equivclass{M}{\approx}}{!A}$ se e somente se $!A \in \Gamma^*$.  Em
  particular, $\Sat{\equivclass{M}{\approx}}{!A}$ para todo $!A \in
  \Gamma$, de modo que $\Gamma$ é satisfatível.}{}
\end{proof}

\begin{cor}[Teorema da Completude, Segunda Versão]
\ollabel{cor:completeness}
Para todo $\Gamma$ e !!{sentence}s~$!A$: se $\Gamma \Entails !A$ então
$\Gamma \Proves !A$.
\end{cor}

\begin{proof}
Note que os $\Gamma$ em \olref{cor:completeness} e
\olref{thm:completeness} são universalmente quantificados.  Para nos
certificarmos de que não nos confundimos, reformulemos \olref{thm:completeness}
usando uma variável diferente: para qualquer conjunto de !!{sentence}s~$\Delta$, se
$\Delta$ é consistente, então é satisfatível.  Por contraposição, se
$\Delta$ não é satisfatível, então $\Delta$ é inconsistente.  Usaremos
isso para provar o corolário.

Suponha que $\Gamma \Entails !A$.  Então $\Gamma \cup \{\lnot !A\}$ é
insatisfatível por \olref[syn][sem]{prop:entails-unsat}.  Tomando $\Gamma
\cup \{\lnot !A\}$ como nosso $\Delta$, a versão anterior de
\olref{thm:completeness} nos dá que $\Gamma \cup \{\lnot !A\}$ é
inconsistente.  Por
\iftag{FOL}{%
  \tagrefs{prfAX/{fol:axd:prv:prop:prov-incons},
    prfSC/{fol:seq:prv:prop:prov-incons},
    prfND/{fol:ntd:prv:prop:prov-incons},
    prfTab/{fol:tab:prv:prop:prov-incons}}}{%
  \tagrefs{prfAX/{pl:axd:prv:prop:prov-incons},
    prfSC/{pl:seq:prv:prop:prov-incons},
    prfND/{pl:ntd:prv:prop:prov-incons},
    prfTab/{pl:tab:prv:prop:prov-incons}}},
$\Gamma \Proves !A$.
\end{proof}

\tagprob{FOL}
\begin{prob}
Use \olref[fol][com][cth]{cor:completeness} para provar
\olref[fol][com][cth]{thm:completeness}, mostrando assim que as duas
formulações do teorema da completude são equivalentes.
\end{prob}
\tagendprob

\tagprob{notFOL}
\begin{prob}
Use \olref[pl][com][cth]{cor:completeness} para provar
\olref[pl][com][cth]{thm:completeness}, mostrando assim que as duas
formulações do teorema da completude são equivalentes.
\end{prob}
\tagendprob

\tagprob{FOL}
\begin{prob}
Para que um sistema de !!{derivation} seja completo, suas regras devem ser
fortes o suficiente para provar que todo conjunto insatisfatível é inconsistente.  Quais das
regras de !!{derivation} foram necessárias para provar a completude?  Alguma
dessas regras não é usada em nenhum lugar da prova?  Para responder
a essas perguntas, faça uma lista ou diagrama que mostre quais das regras
de !!{derivation} foram usadas em quais resultados que levam à prova de
\olref[fol][com][cth]{thm:completeness}.  Certifique-se de anotar quaisquer
usos tácitos de regras nessas provas.
\end{prob}
\tagendprob

\tagprob{notFOL}
\begin{prob}
Para que um sistema de !!{derivation} seja completo, suas regras devem ser
fortes o suficiente para provar que todo conjunto insatisfatível é inconsistente.  Quais das
regras de !!{derivation} foram necessárias para provar a completude?  Alguma
dessas regras não é usada em nenhum lugar da prova?  Para responder
a essas perguntas, faça uma lista ou diagrama que mostre quais das regras
de !!{derivation} foram usadas em quais resultados que levam à prova de
\olref[pl][com][cth]{thm:completeness}.  Certifique-se de anotar quaisquer
usos tácitos de regras nessas provas.
\end{prob}
\tagendprob

\end{document}
